\documentclass[twocolumn]{article}

\usepackage[margin=1.25in]{geometry}
\usepackage{fontspec}
\usepackage{xeCJK}
\usepackage{url} % 加载url包
\usepackage{float}
\usepackage{array}
\geometry{a4paper,left=2cm,right=2cm,top=2.5cm,bottom=2.5cm,columnsep=1cm}
\defaultfontfeatures{Ligatures=TeX,Scale=MatchLowercase}
\setmainfont{Latin Modern Roman}
\setsansfont{Latin Modern Sans}
\setmonofont{Latin Modern Mono}
% 中文主字体：指定 BoldFont 以支持 \textbf 等加粗命令（Windows 常用 SimHei）
% 如果系统无 SimHei，可改为已安装的黑体或思源黑体等字体名
\setCJKmainfont[BoldFont={SimHei}]{SimSun}
\usepackage{authblk}
\usepackage{setspace}
\usepackage{graphicx}
\graphicspath{{./figures/}}
\usepackage{subcaption}
\usepackage{amsmath}
\usepackage{amssymb}
\usepackage{cuted}
\usepackage{booktabs}
\usepackage{tabularx}
\usepackage{float}
\usepackage[section]{placeins}
\usepackage{hyperref}
\usepackage{tikz}
\usetikzlibrary{arrows.meta, positioning, shapes.geometric, calc}
\usepackage{caption}
\usepackage{adjustbox}
\usepackage[vlined]{algorithm2e}
\SetKwInput{KwIn}{Input}
\SetKwInput{KwOut}{Output}

% ---- Column & line number spacing tweaks ----
\setlength{\columnsep}{24pt}      % space between the two columns (default ~10pt)
\usepackage[switch]{lineno}       % 'switch' puts line numbers on outer margins
\setlength\linenumbersep{8pt}     % gap between line numbers and text


%%%%%% Bibliography %%%%%%
% Replace "sample" in the \addbibresource line below with the name of your .bib file.
\usepackage[
  style=nejm, 
  citestyle=numeric-comp,
  sorting=none
]{biblatex}
\addbibresource{sample.bib}

%%%%%% Title %%%%%%
% Full titles can be a maximum of 100 characters, including spaces. 
% Title Format: Use title case, capitalizing the first letter of each word, except for certain small words, such as articles and short prepositions
\title{基于 GDL 的随机纤维网格生成与薄膜桥状 Binder 建模的研究与分析}

%%%%%% Authors %%%%%%
% Authors should be listed in order of contribution to the paper, by first name, then middle initial (if any), followed by last name.
% Authors should be listed in the order in which they will appear in the published version if the manuscript is accepted. 
% Use an asterisk (*) to identify the corresponding author, and be sure to include that person’s e-mail address. Use symbols (in this order: †, ‡, §, ||, ¶, #, ††, ‡‡, etc.) for author notes, such as present addresses, “These authors contributed equally to this work” notations, and similar information.
% You can include group authors, but please include a list of the actual authors (the group members) in the Supplementary Materials.
\author[1*$\dag$]{耿若冰} % 1*$\dag$ 表示的是第一作者，$\dag$ 表示通讯作者，
\author[2]{王晨羽} % 2 表示的是第二作者
\author[2]{刘昕怡} % 3 表示的是第三作者
\author[2]{王学敏} % 4 表示的是第四作者

%%%%%% Affiliations %%%%%%
\affil[1]{卓越 2502，长沙理工大学卓越工程师学院}
\affil[2]{卓越 2501，长沙理工大学卓越工程师学院}

%%%%%% Date %%%%%%
% Date is optional
\date{2026 年 2 月 20 日}

%%%%%% Spacing %%%%%%
% For a proof-style two-column layout, a mild stretch often looks good.
% You can change 1.15 to 1.0 for true single spacing, or back to \onehalfspacing.
\setstretch{1.15}

% 每段首行缩进两格（约为两个中文字符宽度）并启用首段缩进
\usepackage{indentfirst}
\setlength{\parindent}{2em}

% 调整显示公式上下间距，使上下间距一致且较为紧凑
% 这些参数影响由 amsmath 提供的 display math（\[...\], equation 等）
\setlength{\abovedisplayskip}{6pt}
\setlength{\belowdisplayskip}{6pt}
\setlength{\abovedisplayshortskip}{6pt}
\setlength{\belowdisplayshortskip}{6pt}

% 若使用 cuted 的 strip 环境，确保 strip 周围的额外间距也被控制
\makeatletter
\@ifundefined{stripsep}{\newskip\stripsep \stripsep=6pt}{\setlength{\stripsep}{6pt}}
\makeatother

\renewenvironment{abstract}{%
  \par\noindent
  \textbf{Abstract:} % 加粗的“Abstract.”
  \ignorespaces
}{%
}
  \par\vspace{1em} % 摘要与关键词之间的间距

\begin{document}
\twocolumn[
\begin{@twocolumnfalse}

\maketitle

%%%%%% Abstract (two-column) %%%%%%
\noindent\hrule
\vspace{0.5em}
\begin{abstract}
气体扩散层（GDL）作为质子交换膜燃料电池（PEMFC）的核心部件，其微观孔隙结构与纤维排布决定着气、水、电传输效率，开展微观结构建模对提升电池性能具有重要理论与工程价值。本文采用分层随机生成策略构建碳纤维骨架结构：通过高斯抖动、均匀随机采样与双倾角策略，分别确定纤维的中心位置与空间取向，并结合动态加权随机选择方法，使纤维排布呈现 “乱中有序” 的分布特征。针对粘结剂（binder），基于纤维接触与近接触识别算法，分别构建共面胶膜与异面胶桥，以还原其在实际结构中的多样形态。在基础模型生成后，采用高斯平滑处理消除台阶效应与锯齿伪影，提升微观结构的真实感。算法效率层面，引入厚度轴偏移（nudge）与自适应参数收紧等防jamming策略，有效降低结构生成耗时。
\end{abstract}
\\\noindent\textbf{Keywords:} GDL, 随机纤维网络, 薄膜桥状 Binder, 网格生成, 建模，高斯平滑
\vspace{0.5em}
\noindent\hrule
\vspace{1em}
\end{@twocolumnfalse}
]


%%%%%% Main Text %%%%%%

\section{Introduction}



质子交换膜燃料电池(PEMFC)作为高效清洁的能源转换装置，在新能源汽车和分布式发电领域具有广阔的应用前景。气体扩散层(GDL)作为其核心组件，由碳纤维和粘结剂构成，其微观结构直接决定了气体传输、水管理和电池整体性能。因此，精确重构GDL的三维微观结构，对于理解其内部传质机理、优化材料设计具有重要的理论与现实意义。

目前，国内外学者已提出多种GDL微观结构重构方法，主要包括基于实验图像的重构和基于随机算法的生成式重构。然而，现有研究在GDL三维重构中仍存在显著局限：多数随机重构模型不仅将碳纤维简化为理想的长直圆柱体，还进一步假设其在整个厚度方向上均匀分布，完全忽略了实际生产中纤维的交错重叠与局部密度不均(Zhang et al, 2022；Gao et al, 2020)。这种双重简化导致重构的孔隙结构与真实GDL存在明显偏差，进而影响了传质模拟的准确性。同时，X射线CT等实验技术也难以有效区分碳纤维与粘结剂相(Benz et al, 2025；Jin et al, 2019)，导致对其分布的定量分析存在困难。由于上述简化，现有模型往往无法准确预测GDL在压缩等工况下的内部结构变化(Wang et al, 2022)，使得模拟得到的气体渗透、水管理等性能与实验结果存在偏差，限制了模型对GDL设计优化的指导能力。

尽管已有研究尝试通过分层堆叠的方式模拟GDL的厚度方向结构（Lai et al, 2026），但大多仅调控各层的孔隙率，而未对纤维的空间取向与位置分布进行精细化控制。为了在计算可行性与物理真实性之间取得平衡，本研究在现有随机重构模型的基础上，重点突破“均匀分布”的假设，提出了分层取向与分层位置采样策略，对纤维的空间分布与取向进行更精细的建模，并进一步引入粘结剂的分布模型，以更贴近真实GDL的复杂结构。

本研究的目标是：(1)重构出给定孔隙率的碳纤维（重叠/不重叠）气体扩散层三维结构；(2)考虑粘结剂(binder)的影响，重构出更贴近真实情况的GDL三维结构。通过本研究，旨在提供一套完整的、可执行的三维重构流程与代码，并生成相应的三维数据文件，为后续的传质模拟和性能分析提供可靠的基础模型。



%%%%%% Equations %%%%%%
% \subsection*{Equations}
% Equations should be provided in a text format, rather than as an image. Equations should be numbered consecutively, in round brackets, on the right-hand side of the page by using the ``\textbackslash begin\{equation\}'' command. They should be referred to as Equation 1, etc. in the main text.

% \medskip For example, see Equation \ref{eq:1} and Equation \ref{eq:2} below.
% \begin{equation} \label{eq:1}
%     a^2 + b^2 = c^2
% \end{equation}
% \begin{equation} \label{eq:2}
% \begin{split}
% A & = \frac{\pi r^2}{2} \\
%  & = \frac{1}{2} \pi r^2
% \end{split}
% \end{equation}

\section{Methods}

\subsection{\textbf{计算域与问题定义}}

为了能够在后续的建模和仿真中准确描述碳棒纤维与粘连剂（binder）的几何关系，本节首先给出计算域、变量定义与基本假设。我们在尺寸为 $I \times J \times K$ 的三维长方体离散体素域（3D voxel domain）$\Omega \subset \mathbb{Z}^3$ 中生成纤维（碳棒）多孔介质，$I,J,K$ 的具体值由实验样品的实际尺寸决定，除非特别说明，本文均以 $155 \times 300 \times 300$ 为例。

\subsubsection{\textbf{纤维几何表示与体素化}}

每根纤维被建模为一个圆柱体，具有固定的半径 $r$，其轴向长度 $l \in [l_{\min}, l_{\max}]$。本文用一条三维线段 $[\textbf a, \textbf b]$ 来表示每根纤维的中心轴线，其中 $\textbf a,\textbf b \in \mathbb{R}^3$ 分别表示中心轴线起点和终点的三维坐标。纤维固相体素集合记为 $V_{\textbf{fiber}}$。对任意体素索引 $\textbf{x} = (i,j,k) \in \Omega$，定义状态函数 $S(\textbf x)$ 来表示该体素是否被纤维占据：
\[
S(\textbf x) = 
\begin{cases}
1, & \textbf{x} \in V_{\textbf{fiber}}  \\
0, & \textbf{x} \notin V_{\textbf{fiber}} 
\end{cases}
\]
设计算域内体素总数为
\[
N = I \times J \times K
\]
当前纤维固相体素数为
\[
N_s = \sum \limits_{\textbf{x} \in \Omega} S(\textbf x) = \sum \limits_{i,j,k} S(i,j,k)
\]

\subsubsection{\textbf{binder 与孔隙相的定义}}

为描述粘连剂分布，本文进一步定义 binder 体素集合 $V_{\textbf{binder}} \in \Omega$，以及对应的状态函数 $B(\textbf x)$：
\[
B(\textbf x) =
\begin{cases}
1, & \textbf{x} \in V_{\textbf{binder}}  \\
0, & \textbf{x} \notin V_{\textbf{binder}} 
\end{cases}
\]
相应地，孔隙体素集合定义为
\[
V_{\textbf{pore}} = \Omega \setminus (V_{\textbf{fiber}} \cup V_{\textbf{binder}}) = (V_{\textbf{fiber}} \cup V_{\textbf{binder}})^c
\]
其中 $(\cdot)^c$ 表示相对于 $\Omega$ 的补集，对应的状态函数为
\[
P(\textbf x) = 1 - S(\textbf x) - B(\textbf x)
\]
本文默认三相体素互斥，即 $V_{\textbf{fiber}} \cap V_{\textbf{binder}} = \emptyset$。据此，纤维体积分数 $\phi_{f}$、binder 体积分数 $\phi_{b}$ 和孔隙率 $\varepsilon$ 分别定义为：
\[
\begin{cases}
\phi_{f} = \frac{1}{N} \sum \limits_{\textbf{x} \in \Omega} S(\textbf x) \\
\phi_{b} = \frac{1}{N} \sum \limits_{\textbf{x} \in \Omega} B(\textbf x) \\
\varepsilon = \frac{1}{N} \sum \limits_{\textbf{x} \in \Omega} P(\textbf x) = 1 - \phi_{f} - \phi_{b}
\end{cases}
\]
当不考虑 binder 时，有 $\phi_{b} = 0$，于是孔隙率退化为 $\varepsilon = 1 - \phi_{f}$。

\subsubsection{\textbf{建模目标与约束}}

本文设定目标孔隙率为 $\varepsilon_{\textbf{target}}$，并（可选地）设定目标 binder 体积分数 $\phi_{b,\textbf{target}}$，以指导后续的纤维与 binder 生成过程。对于碳棒纤维，相应地，总固相体积分数满足
\[
\phi_{\textbf{solid,target}} = 1 - \varepsilon_{\textbf{target}} = \phi_{f} + \phi_{b}
\]
其对应地目标固相体素数为
\[
N_{\textbf{solid,target}} = \phi_{\textbf{solid,target}} \cdot N
\]
本文的建模目标是通过逐根生成碳棒纤维（以及随后生成 binder），同时考虑允许重叠和不重叠两种情况，使得最终孔隙率满足
\[
|\varepsilon - \varepsilon_{\textbf{target}}| \le \delta_{\varepsilon}
\]
等价地，使得总固相体素数逼近目标值
\[
N_s^{(\textbf{total})} = \sum \limits_{\textbf{x} \in \Omega} (S(\textbf x) + B(\textbf x)) \approx N_{\textbf{solid,target}}
\]
其中 $\delta_{\varepsilon}$ 是允许误差阈值。后续章节将分别给出重叠和不重叠两种情况下纤维几何的生成以及 binder 的生成的具体算法与实现细节。

\subsection{\textbf{碳棒纤维生成策略}}

\subsubsection{\textbf{算法概述、流程与关键参数}}

本节在 2.1 给出的计算域 $\Omega$、状态函数 $S(\textbf x)$、孔隙率目标约束 $|\varepsilon - \varepsilon_{\textbf{target}}| \le \delta_{\varepsilon}$ 的基础上，构建一种面向 GDL 的分层随机纤维生成算法，用于生成纤维相体素集合 $V_{\textbf{fiber}}$。算法以“逐根生成——判定——写入”的迭代方式进行：每次迭代采样候选纤维的中心位置，方向与长度并进行边界裁剪；允许重叠时直接体素化写入，禁止重叠时先进行碰撞检测，通过后写入；每次写入后更新当前孔隙率 $\varepsilon$，当满足目标约束或达到最大尝试次数时终止。
具体的算法流程如下所示：

% 将流程图以内嵌方式显示（非浮动），避免被 LaTeX 浮动机制推到别处或造成两栏不平衡的大空白 
\begin{center} % 将流程图限制在单栏宽度内并在必要时缩放 
  \begin{adjustbox}{max width=\columnwidth,center} 
    \begin{tikzpicture}[ 
      node distance=5mm and 10mm,
      font=\scriptsize, 
      startstop/.style={rectangle, rounded corners, draw, fill=white, align=center, text width=26mm, inner sep=3pt, minimum height=6mm}, 
      process/.style={rectangle, draw, fill=white, align=center, text width=30mm, inner sep=3pt, minimum height=6mm}, decision/.style={diamond, aspect=2.1, draw, fill=white, align=center, text width=28mm, inner sep=1.5pt}, arrow/.style={-Latex, line width=0.6pt} ] 
      \node (start) [startstop] {开始}; 
      \node (init) [process, below=of start] {初始化：$S(x)=0$，纤维列表为空\\计算当前孔隙率 $\varepsilon$}; 
      \node (stopq) [decision, below=of init, yshift=-2mm] {是否满足停止条件？\\$|\varepsilon-\varepsilon_{\text{target}}|\le\delta_\varepsilon$\\或 tries 达上限}; 
      \node (out) [startstop, right=of stopq, xshift=12mm] {输出：$S(x)$ 与纤维列表}; 
      \node (layer) [process, below=of stopq, yshift=-2mm] {选择层 $m$ 并采样中心 $c$\\（分层 + jitter）}; 
      \node (ori) [process, below=of layer] {采样取向 $(\alpha,\beta)$\\得到方向向量 $u$}; 
      \node (len) [process, below=of ori] {采样长度 $l\in[l_{\min},l_{\max}]$\\构造端点 $a,b$}; 
      \node (clip) [process, below=of len] {边界裁剪（clipped）\\得到域内线段 $[a,b]$}; 
      \node (over) [decision, below=of clip, yshift=-2mm] {是否允许重叠？}; 
      \node (vox) [process, left=of over, xshift=-12mm] {体素化写入 $S(x)$\\记录纤维信息}; 
      \node (coll) [process, right=of over, xshift=12mm] {碰撞检测（不重叠）\\分箱/AABB 预筛\\线段最小距离 $\ge 2r$}; 
      \node (pass) [decision, below=of coll, yshift=-2mm] {检测通过？}; 
      \node (adapt) [process, below=of pass] {失败计数+1\\局部 nudge / 自适应缩短 $l_{\max}$\\缓解 jamming}; 
      \node (upd) [process, above=of vox] {更新孔隙率 $\varepsilon$}; 
      \draw [arrow] (start) -- (init); 
      \draw [arrow] (init) -- (stopq); 
      \draw [arrow] (stopq) -- node[above] {是} (out); 
      \draw [arrow] (stopq) -- node[left] {否} (layer); 
      \draw [arrow] (layer) -- (ori); 
      \draw [arrow] (ori) -- (len); 
      \draw [arrow] (len) -- (clip); 
      \draw [arrow] (clip) -- (over); 
      \draw [arrow] (over) -- node[above] {允许} (vox); 
      \draw [arrow] (over) -- node[above] {不允许} (coll); 
      \draw [arrow] (coll) -- (pass); 
      % 从 pass.west 出发：先向左，再向上，汇入 over->vox 的“允许”箭头（保持黑色、无箭头）
      \coordinate (pass_left) at ($(pass.west)+(-72mm,0)$);
      \coordinate (allow_mid) at ($(over.west)!0.5!(vox.east)$);
      \coordinate (pass_left_up) at ($(pass_left |- allow_mid)$);
      \draw [line width=0.6pt] (pass.west) -- node[below,pos=0.15]{是} (pass_left) -- (pass_left_up) -- (allow_mid);
      \draw [arrow] (pass) -- node[left] {否} (adapt); 
      % 从 adapt.east 出发：先向右一段，再向上对齐 stopq.south 的纵坐标，最后水平连到 stopq.south（改为右侧回路）
      \coordinate (adapt_right) at ($(adapt.east)+(12mm,0)$);
      \coordinate (adapt_right_up) at ($(adapt_right |- stopq.south)$);
      \draw [arrow] (adapt.east) -- (adapt_right) -- (adapt_right_up) -- (stopq.south);
      \draw [arrow] (vox) -- (upd); 
      \draw [arrow] (upd.north) |- (stopq.west); 
    \end{tikzpicture} 
  \end{adjustbox} 
  \captionof{figure}{分层随机纤维生成流程（含允许/不允许重叠分支与自适应防 jamming 策略）。} 
  \label{fig:fiber_flowchart} 
\end{center}
分层随机纤维生成算法的主要参数如下：

\begin{table}[htbp]
\centering
\begin{adjustbox}{max width=\columnwidth,center}
\begin{tabular}{lll}
\hline
\textbf{符号/名称} & \textbf{含义} & \textbf{典型取值/范围} \\
\hline
$I,J,K$ & 体素域尺寸（第 2.1.1 节） & $155\times 300\times 300$ \\
$r$ & 纤维半径（体素） & 例如 $r=3.5$ \\
$l_{\min},l_{\max}$ & 纤维长度范围（体素） & 例如 $[140,320]$ \\
$n_{\mathrm{layer}}$ & 厚度方向分层数 & 例如 $6\sim 11$ \\
$\sigma_z$（jitter） & 层内厚度方向中心抖动 & 例如 $2\sim 5$ 体素 \\
$\mu_{\alpha}^{(m)}$ & 第 $m$ 层的平面内平均主方向 & 给定/缓慢变化 \\
$\sigma_\alpha$ & 平面内角度离散度（层内“凌乱”程度） & 例如 $40^\circ\sim 80^\circ$ \\
$\sigma_\beta,\beta_{\max}$ & 离层倾角控制（贴层程度） & 例如 $\sigma_\beta=2^\circ\sim 5^\circ$ \\
\texttt{overlap} & 是否允许重叠（Y/N） & Y / N \\
\texttt{fiber\_end\_mode} & 端点边界处理方式 & clipped \\
$\varepsilon_{\mathrm{target}}$ & 目标孔隙率（第 2.1.3 节） & 例如 $0.70\sim 0.85$ \\
$\delta_\varepsilon$ & 孔隙率容差（停止阈值） & 例如 $10^{-4}\sim 10^{-3}$ \\
\texttt{max\_tries} & 最大尝试次数（防止无限循环） & 依实现设定 \\
\hline
\end{tabular}
\end{adjustbox}
\caption{分层随机纤维生成算法的主要参数。}\label{tab:fiber_params_cn}
\end{table}

\subsubsection{\textbf{厚度方向分层建模与层内位置采样}}

GDL 在厚度方向通常呈现\textbf{厚度方向层状堆叠}、\textbf{层内位置随机、乱中有序}的特征，本节主要通过分层建模与层内位置采样，来给每根纤维分配一个合理的厚度位置，并给它采样一个中心点 $\textbf c = (x,y,z)$。

取计算域中最短的维度（通常是 $I$ 轴，对应实现中的 TP\_AXIS）作为厚度方向（through-plane），其余两维为面内方向（in-plane）（对应实现中的 IP1\_AXIS 和 IP2\_AXIS）。设厚度方向的长度为 $K_{\textbf{TP}}$，为贴近实际 GDL 的结构，我们将厚度方向划分为 $n_{\mathrm{layer}}$ 个连续的层（layer），那么每一层的厚度（层厚）为
\[
h = \dfrac{K_{\textbf{TP}}}{n_{\mathrm{layer}}}
\]
其中第 $m$ 层的中心位置为
\[
z_m = \left(m + \dfrac 1 2\right)h, \quad m=0,1,\ldots,n_{\mathrm{layer}}-1
\]

考虑如何对于每根碳棒纤维选择一个层。对于\textbf{允许重叠}的情况，我们直接等概率的选取一个层 $m$；而对于\textbf{禁止重叠}的情况，为了缓解后续可能出现的 jamming（即后续纤维难以找到合适位置写入），我们采用一种基于当前层占用情况的\textbf{动态加权随机选择}策略：为了避免过早地将某些层填满导致后续纤维难以写入，我们应当让当前占用较少的层有更高的概率被选中，具体地，我们定义
\[
w_i = \dfrac{1}{cnt_i + 1}, \quad i=0,1,\ldots,n_{\mathrm{layer}}-1
\]
其中 $cnt_i$ 表示第 $i$ 层当前已写入的纤维数量，$w_i$ 则是第 $i$ 层被选中的权重。
通过这种方式，不仅可以在禁止重叠的情况下更均匀地分布纤维，避免过早地填满某些层导致后续纤维难以写入，还可以保留一定的随机性，使得生成的结构更接近实际 GDL 的“乱中有序”特征。

在前文中，我们已经找到了每一层的中心位置 $z_m$，接下来我们需要为每根纤维采样一个具体的厚度位置 $z$。为了模拟实际 GDL 中纤维在层中心较为集中、向上下逐渐稀疏的分布特征，我们在层中心位置 $z_m$ 的基础上，添加一个服从零均值正态分布的高斯抖动（jitter）：
\[
z \sim \mathcal{N}(z_m,\sigma^2)
\]
其中 $z$ 是最终的厚度位置，$\sigma$ 作为标准差则控制了层内厚度方向的离散程度（即“凌乱”程度）。根据正态分布的特点，最终大多数点会落在 $z_m$ 附近，距离中心越远概率快速变小；较小的 $\sigma$ 会使得纤维更集中在层中心，而较大的 $\sigma$ 则会使得纤维在厚度方向上分布得更为分散。通过调整 $\sigma$ 的值，我们可以模拟不同程度的层内“凌乱”，以更好地匹配实际 GDL 的结构特征。

需要注意的是，正态分布的支持是整个实数轴，因此理论上可能会采样到超出计算域厚度范围的 $z$ 值。为了确保生成的纤维能够正确地写入计算域，我们需要对采样得到的 $z$ 进行边界裁剪（clipping），将其限制在合理的范围内。

接下来，我们需要为每根纤维采样一个面内位置 $(x,y)$，为了满足 GDL 层内“凌乱”的特性，我们可以直接在计算域的 $xy$ 平面上进行均匀随机采样：
\[  
\begin{cases}
x \sim \mathcal{U}(0, I_{\textbf{IP1}}) \\
y \sim \mathcal{U}(0, I_{\textbf{IP2}})
\end{cases}
\]
其中 $I_{\textbf{IP1}}$ 和 $I_{\textbf{IP2}}$ 分别表示计算域在面内两个方向上的尺寸。通过这种方式，能够确保每一层内的纤维在面内位置上是随机分布的。据此，我们就得到了每根纤维的中心点坐标 $\textbf c = (x,y,z)$。

通过上述算法，我们成功地为每根纤维分配了一个合理的中心点坐标，并且模拟了实际 GDL 中纤维在厚度方向层状堆叠、层内位置随机但又有一定规律的分布特征。这为后续的纤维生成与写入奠定了坚实的基础。

\subsubsection{\textbf{纤维方向采样与空间位置确定}}

在上一节的内容中，我们已经为每根纤维分配了一个合理的中心点坐标。那么为了唯一确定一根三维空间中的碳棒纤维，我们还需要指定它的方向。

我们用两个角度来参数化纤维的方向：
\begin{itemize}
\item $\alpha$：纤维在面内（in-plane azimuth）的主方向角，定义为纤维轴线在 $xy$ 平面上的投影与 $x$ 轴的夹角，取值范围为 $\alpha \in [0,2 \pi)$。
\item $\beta$：纤维的离层倾角（out-of-plane tilt），定义为纤维轴线与厚度方向（$z$ 轴）的夹角。
\end{itemize}
则纤维的方向向量可定义为：
\[
\textbf u = (\cos\alpha \cos\beta, \cos\beta \sin\alpha, \sin\beta)
\]

记第 $m$ 层的平均主方向为 $\mu_{\alpha}^{(m)}$，我们在层内采样 $\alpha$ 时，以 $\mu_{\alpha}^{(m)}$ 为均值，使用一个适当的标准差 $\sigma_\alpha$ 来控制层内纤维方向的离散程度（即“凌乱”程度），从而使得每层存在一个明显的以 $\mu_{\alpha}^{(m)}$ 为主方向，纤维倾向于围绕它分布的特征。具体地，我们可以使用一个截断高斯 + 环上高斯来采样 $\alpha$：
首先采样一个“未取模”的角度：
\[
\tilde\alpha \sim \mathcal{N}(\mu_{\alpha}^{(m)}, \sigma_\alpha^2)
\]
再把它映射到 $[0,2\pi)$ 上：
\[
\alpha = \tilde\alpha \bmod 2\pi
\]
同时，为了防止高斯长尾导致过于离谱的角度，规定最大偏转角 $\alpha_{\max}$，在采样 $\tilde\alpha$ 时进行截断，限制其在 $[\mu_{\alpha}^{(m)} - \alpha_{\max}, \mu_{\alpha}^{(m)} + \alpha_{\max}]$ 范围内。
对于 $\mu_{\alpha}^{(m)}$ 的设置，可以让所有层共享一个全局主方向 $\mu_{\alpha}$；也可以让每层主方向 $\mu_{\alpha}^{(m)}$ 随层数 $m$ 缓慢变化，即 $\mu_{\alpha}^{(m)} = \mu_0 + \Delta\mu \cdot m$，以模拟实际 GDL 中不同层主方向逐渐变化的特征。

对于离层倾角 $\beta$ 的采样，我们同样使用一个截断高斯来控制纤维的贴层程度（即与厚度方向的夹角）。由于 GDL 的真实结构中纤维大多“躺在层平面内”，所以设定 $\beta$ 的均值为 $0^\circ$，表示纤维倾向于贴近层面。即 $\beta$ 满足正态分布
\[
\tlide\beta \sim \mathcal{N}(0, \sigma_\beta^2), \quad \beta = \text{clip}(\tilde\beta, -\beta_{\max}, \beta_{\max})
\]
其中 $\beta_{\max}$ 则限制了纤维的最大离层倾角，相当于是一个硬性约束。假设一根纤维的中心轴线长度为 $L$，与层平面的夹角为 $\beta$，那么应当满足
\[
L \sin|\beta| \le K_{\textbf{TP}}
\]
因为对任意 $L \ge L_{\min}$，所以要使得上式满足对于所有可能的纤维长度，则需要满足
\[
L_{\min} \sin\beta \le K_{\textbf{TP}}
\]
化简得
\[
|\beta| \le \sin^{-1}\left(\dfrac{K_{\textbf{TP}}}{L_{\min}}\right)
\]
于是定义
\[
\beta_{\max} = \sin^{-1}\left(\dfrac{K_{\textbf{TP}}}{L_{\min}}\right)
\]

由上述分布得到的方向向量 $\textbf u$，结合前面采样得到的中心点 $\textbf c$，就可以唯一确定一根纤维的空间位置和方向。

\subsubsection{\textbf{纤维长度与端点处理}}

在前面的内容中，我们已经通过中心点 $\textbf c$ 和方向向量 $\textbf u$ 确定了纤维的空间位置和方向。接下来，我们需要为每根纤维采样一个长度 $l$，以构造出完整的纤维线段 $[\textbf a, \textbf b]$，其中 $\textbf a$ 和 $\textbf b$ 分别是纤维的起点和终点坐标。

纤维的长度 $l$ 直接影响着纤维在空间中的占用范围以及与其他纤维的交互情况。为了模拟实际 GDL 中纤维长度的分布特征，我们通常设定一个合理的长度范围 $[l_{\min}, l_{\max}]$，并从中采样一个长度值 $l$。具体地，我们可以使用一个均匀分布来采样长度：
\[
l \sim \mathcal{U}(l_{\min}, l_{\max})
\]
需要注意的是，在不重叠模式下，较长的纤维更容易与已有纤维发生碰撞，因此在实际实现中，我们可以根据当前的填充情况动态调整 $l_{\max}$，以缓解 jamming 问题。

对于每根纤维，在已知长度 $l$，中心点坐标 $\textbf c$ 和方向向量 $\textbf u$ 时，我们可以计算出纤维的起点 $\textbf a$ 和终点 $\textbf b$：
\[
\begin{cases}
\textbf a = \textbf c - \dfrac{l}{2} \textbf u \\
\\
\textbf b = \textbf c + \dfrac{l}{2} \textbf u
\end{cases}
\]
上述构造保证线段长度为 $|b-a|=l$，且线段中点为 $c$。该线段 $[a,b]$ 为候选中心轴线，后续将根据计算域边界执行裁剪（clipped）以得到位于域内的有效线段，详见第 2.2.5 节。

\subsubsection{\textbf{边界裁剪与体素化写入}}




\subsection{\textbf{binder 生成策略}}

\subsubsection{\textbf{算法概述、输入输出与关键参数}}

在纤维骨架生成完成后（纤维相几何分布 $S(\textbf x)$ 已固定且不再发生改变）。在此基础上，本文进一步开展 binder 相 $B(\textbf x)$ 的重构于生成，以表征粘连剂在纤维接触/近接触区域富集形成胶桥/胶膜的微结构特征。生成过程中，通过控制粘结剂的沉积位置、分布形态与填充范围，在保证物理合理性的前提下，使最终 binder 体积分数 $\phi_{b}$ 以及总体孔隙率 $\varepsilon$ 收敛至预设的目标值 $\phi_{b,\textbf{target}}$，从而实现对气体扩散层多相微结构的可控、定量重建，为后续输运性能分析提供符合真实材料特征的几何模型。

总体上，binder 生成遵循“\emph{接触识别-形态分类-局部几何构造-体积分数校准}”的流程：首先在纤维几何集合中识别近接触纤维对并确定最近点对作为锚点；随后根据两纤维相对取向（共面/异面）构造不同形态的三维局部 binder 体积；将其体素化写入 $B(\textbf x)$ 并保持与纤维互斥（$B(\textbf x) \cdot S(\textbf x) = 0$）；最后通过调整局部几何或迭代生成，使 $\phi_{b}$ 逼近目标 $\phi_{b,\textbf{target}}$。

\noindent\mbox{流程摘要如下：}
\begin{algorithm}[t]
\small
\DontPrintSemicolon
\SetAlgoVlined
\SetKw{Break}{break}
\SetKwFor{ForEach}{for each}{do}{end}

\BlankLine
\hrule
\BlankLine

\KwIn{$S\in\{0,1\}^{\Omega}$；$\mathcal{F}=\{(a_n,b_n,u_n,\ell_n)\}_{n=1}^{N_f}$；$\phi_{b}^{\star}$；$d_c$；$\theta_0$；$\Theta$}
\KwOut{$B\in\{0,1\}^{\Omega}$，with $B\odot S=0$}

$B \leftarrow 0$;\quad $N_b^{\star}\leftarrow \phi_b^{\star}\,|\Omega|$;\quad $N_b\leftarrow 0$\;
$\mathcal{C}\leftarrow \{(i,j): d([a_i,b_i],[a_j,b_j])\le d_c\}$，并记录 $(\textbf A_{ij},\textbf B_{ij})$\;

\ForEach{$(i,j)\in\mathcal{C}$}{
    $\theta_{ij}\leftarrow \arccos(u_i\cdot u_j)$;\quad
    $T_{ij}\leftarrow \mathbf{1}[\theta_{ij}\le \theta_0]$ \tcp*{$T_{ij}=1$ 共面，否则异面}
    $\mathcal{V}_{ij}\leftarrow \mathrm{Bridge}(\textbf A_{ij},\textbf B_{ij},T_{ij};\Theta)$\;
    $B \leftarrow B \,\vee\, \mathrm{Voxelize}(\mathcal{V}_{ij})$\;
    $B \leftarrow B \wedge \neg S$\;
    $N_b \leftarrow \sum_{x\in\Omega} B(x)$\;
    \If{$N_b \ge N_b^{\star}$}{\Break}
}

\If{$N_b > N_b^{\star}$}{
    $B \leftarrow \mathrm{Trim}(B,\,N_b^{\star})$ \tcp*{删除多余体素}
}
\Return{$B$}\;

\BlankLine
\hrule

\caption{基于纤维接触几何的 binder 生成流程}
\label{alg:binder}
\end{algorithm}

binder体素场 $S(x)$ 与纤维几何集合 $\mathcal{F}$ 为输入，给定接触阈值 $d_c$、形态分类阈值 $\theta_0$ 以及几何参数 $\Theta$，输出满足互斥约束 $B(x)\cdot S(x)=0$binder der 体素场 $B(x)$，并使 $\phi_b$ 接近目标 $\phi_b^\star$。


\subsubsection{\textbf{纤维接触/近接触识别}}

在气体扩散层（GDL）中，粘结剂（binder）是一种基于碳的固相物质，它通过在纤维间形成结构键维持材料结构完整性，其在碳纸制造中由可碳化热固性树脂经碳化形成，主要分布在纤维接触/近接触处的高表面积区域，为了拟合真实的 binder 分布特征，首先要从纤维集合里找到候选接触对。

记第 $i,j$ 根纤维的中心轴线分别为 $[a_i,b_i]$ 和 $[a_j,b_j]$，定义两线段的最小距离为
\[
\begin{aligned}
d_{i,j} & = d([a_i,b_i],[a_j,b_j]) \\
& = \min_{\textbf A_{ij}\in[a_i,b_i], \textbf B_{ij}\in[a_j,b_j]} \|\textbf A_{ij}-\textbf B_{ij}\|
\end{aligned}
\]
当 $d_{i,j} \le d_c$ 时，认为两根纤维满足接触/近接触条件，即
\[
(i,j)\in\mathcal{C} \quad \Longleftrightarrow \quad d_{ij}\le d_c.
\]
其中 $\mathcal{C}$ 表示候选接触对集合。除距离判定外，本文同时记录实现最小距离的最近点对 $(\textbf A_{ij},\textbf B_{ij})$，其中 $\textbf A_{ij}\in[a_i,b_i]$，$\textbf B_{ij}\in[a_j,b_j]$，用于后续 binder 的局部几何构造。

考虑如何选取合适的阈值 $d_c$。本文将纤维建模为半径为 $r$ 的圆柱体，因此当 $d_{ij} \le 2r$ 时，两根纤维存在实际接触或重叠。由于粘连剂具有一定跨越能力，本文允许在小间隙内形成胶桥，将近接触阈值取为
\[
d_c = 2r + \delta
\]
其中 $\delta \ge 0$ 表示容许间隙参数，用于控制 binder 生成的“跨缝”尺度。

直接对所有纤维对进行距离计算的复杂度为 $O({N_f}^2)$。为降低计算成本，本文在实现中结合厚度分箱与 AABB 预筛对候选对进行快速过滤，仅对潜在相邻纤维对执行精确的线段最短距离计算。

\subsubsection{\textbf{binder 形态分类：共面胶膜与异面胶桥}}

在实际的 GDL 中，binder 在不同几何接触情形下呈现出不同的宏观形态：当两根纤维在同一层平面内近似共面时，binder 更倾向沿接触区域铺展形成胶膜状结构；当两根纤维以较大夹角交叉或跨层接近时，binder 更倾向在局部形成胶桥状结构。为了在生成式模型中体现上述差异，本文对候选接触对进行形态分类，并在后续几何构造中采用不同参数化描述。

我们考虑基于两根纤维的方向夹角作为分类指标。对于前文中得到的候选接触对 $(i,j)\in\mathcal{C}$，记两根纤维的单位方向向量为 $\textbf u_i$ 和 $\textbf u_j$，则它们的夹角 $\theta_{ij}$ 定义为
\[
\theta_{ij} = \arccos\!\left(\left|\textbf u_i\cdot \textbf u_j\right|\right)\in\left[0,\dfrac{\pi}{2}\right]
\]
其中绝对值用于消除方向符号（$\textbf u$ 和 $-\textbf u$ 代表同一纤维方向）。

设定一个夹角阈值 $\theta_0$，我们将接触对分为两类：
\[
T_{ij} = 
\begin{cases}
\text{coplanar}, & \theta_{ij}\le \theta_0,\\
\text{skew}, & \theta_{ij}>\theta_0.
\end{cases}
\]
其中 \texttt{coplanar} 表示共面胶模型，\texttt{skew} 表示异面桥模型。若实现中记录了纤维所属层 $m_i$，则也可结合层信息进行分类，例如仅当 $|m_i - m_j| \le 1$ 时才视为共面型，以增强对跨层接触的鲁棒性。

通过以上分类，我们在后续的 binder 生成过程中针对不同类型的接触对采用不同的几何构造参数：共面型采用更大的沿纤维延伸尺度与更扁平的弯月面厚度分布，以模拟胶膜状结构；异面型采用更局部的延伸尺度与更强的弯月面曲率，以模拟胶桥状结构。该分类使 binder 形态从共面胶膜到异面胶桥可连续过渡，避免采用完全割裂的两套模型。

\subsubsection{\textbf{局部几何构造与体素化写入}}

上文中提到，接触对 $(i,j)\in\mathcal{C}$ 的最近点对为 $(\textbf A_{ij},\textbf B_{ij})$，我们以此为锚点，定义粘连剂的桥接方向：
\[
\textbf{n}_{ij} = \dfrac{\textbf B_{ij}-\textbf A_{ij}}{\|\textbf B_{ij}-\textbf A_{ij}\|}
\]
该方向刻画两纤维间的最近连接法昂想，为后续局部胶桥构造提供参考。

考虑对于每个接触对 $(i,j)$ 如何构建几何上合理的 binder 生成的范围。我们首先选择一个轴向延伸半长 $\ell$（span），对共面/异面可不同：
\[
\ell = 
\begin{cases}
\ell_{\text{coplanar}}, & T_{ij}=\text{coplanar},\\
\ell_{\text{skew}}, & T_{ij}=\text{skew}.
\end{cases}
\]
一般来说，共面型通常采用较大的 $\ell$ 以模拟胶膜沿接触区域铺展，异面型采用较小的 $\ell$ 以模拟局部胶桥聚集。以 $\textbf A_{ij}$ 和 $\textbf B_{ij}$ 为中心点，沿两根纤维轴向对称方向分别延伸 $\ell$，得到四个角点
\[
\begin{cases}
    \textbf P_1 = \textbf A_{ij} - \ell \textbf u_i \\
    \textbf P_2 = \textbf A_{ij} + \ell \textbf u_i,\\
    \textbf Q_1 = \textbf B_{ij} - \ell \textbf u_j \\
    \textbf Q_2 = \textbf B_{ij} + \ell \textbf u_j
\end{cases}
\]
这四点构成了一个扭曲四边形“胶桥骨架”，即一个二维上的桥接中心曲面。

为在两根纤维之间构造连续的局部 binder 几何，本文采用参数化曲面（patch）描述桥接区域的中间面。引入二维参数 $(s,t) \in [0,1]^2$，并定义映射
\[
G:[0,1]^2\rightarrow \mathbb{R}^3,\qquad (s,t)\mapsto G(s,t),
\]
其中 $G(s,t)$ 给出桥接面片在三维空间中的坐标点。假设 $\textbf P_1,\textbf P_2,\textbf Q_1,\textbf Q_2 \in \mathbb{R}^3$，本文采用\textbf{双线性插值（bilinear patch）}\cite{赵煌2012双线性插值算法的优化及其应用}构造桥接面片：
\[
G(s,t) = (1 - s)(1 - t)\textbf P_1 + s(1 - t)\textbf P_2 + (1 - s)t \textbf Q_1 + st \textbf Q_2
\]
该构造满足：
\[
\begin{cases}
G(0,0) = \textbf P_1\\
G(1,0) = \textbf P_2\\
G(0,1) = \textbf Q_1\\
G(1,1) = \textbf Q_2
\end{cases}
\]
从而将单位参数域的四个角与三维空间中的四个角点一一对应，构建了三维空间内 binder 分布的连续曲面描述。

上述 $G(s,t)$ 定义了一个理想化的桥接面片，但实际 binder 具有一定厚度和体积。为此，本文在参数域 $[0,1]^2$ 上定义一个厚度函数：
\[
R:[0,1]^2\rightarrow \mathbb{R}_{+},\qquad (s,t)\mapsto R(s,t),
\]
其中 $R(s,t)$ 表示 $G(s,t)$ 处 binder 的局部“加厚尺度”。直观上，可将 binder 体积视为对面片 $G(s,t)$ 的管状加厚（tubular thickening）：在每个参数点 $(s,t)$ 处，以 $G(s,t)$ 为中心生成半径为 $R(s,t)$ 的局部邻域并取其并集，从而得到三维 binder 体积
\[
\mathcal{V}=\bigcup_{(s,t)\in[0,1]^2}\left\{x\in\mathbb{R}^3:\ \|x-G(s,t)\|_2\le R(s,t)\right\}.
\]
也就是把桥接面片 $G(s,t)$ 上的每一个点，当作一个球心，生成一个半径为 $R(s,t)$ 的球体，然后将这些球体的并集作为 binder 的三维几何表示。我们定义厚度函数的计算方式为：
\[
R(s,t) = R_0 + R_{\textbf{mid}} \sin(\pi s) \sin(\pi t) - R_{\textbf{conc}} \cdot f(s,t)
\]
其中 $R_0$ 表示基础厚度，保证桥接区不出现“断裂/空洞”；$R_{\textbf{mid}}$ 表示沿桥接面片中心区域的加厚幅度，通过乘上双正弦函数来精确模拟 binder 在接触区域的富集特征；$R_{\textbf{conc}}$ 表示对桥接面片边缘的凹弧控制，模拟 binder 在边缘区域的收缩特征。为了得到合理的 binder 凹弧形态，我们定义一个边缘约束函数 $f(s,t)$：
\[
f(s,t) = \sin(\pi s) \sin(\pi t)
\]
则整个函数变为
\[
R(s,t) = R_0 + (R_{\textbf{mid}} - R_{\textbf{conc}}) \sin(\pi s) \sin(\pi t)
\]
通过调整 $R_{\textbf{mid}}$ 和 $R_{\textbf{conc}}$ 的相对大小，可以控制 binder 在中心区域的加厚程度与在边缘区域的凹弧程度，从而模拟出更符合实际 GDL 中 binder 形态的桥接结构。

得到 binder 的三维几何表示 $\mathcal{V}$ 后，接下来需要将其体素化写入 binder 体素场 $B(\textbf x)$ 中。记对接触对 $(i,j)$ 构造的 binder 体积为 $\mathcal{V}_{ij}$，对体素中心点 $\textbf p_x$，若 $\textbf p_x \in \mathcal{V}_{ij}$ 则将 $B(\textbf x)$ 置为 1。为保证三相互斥，写入时加入以下强制约束
\[
B(\textbf x) \leftarrow B(\textbf x) \cdot (1 - S(\textbf x))
\]
即 binder 不占据任何已被纤维占据的体素。同时为了降低计算成本，体素化仅在 $\mathcal{V}_{ij}$ 的轴对齐包围盒（AABB）范围内进行，避免对整个计算域进行无效遍历。

至此，我们完成了基于纤维接触几何的 binder 生成的核心流程：从接触识别到形态分类，再到局部几何构造与体素化写入，最终得到满足预设 binder 体积分数目标的 binder 体素场 $B(\textbf x)$。

\subsubsection{\textbf{体积分数匹配与停止准则}}

binder 生成以体积分数约束为目标，其体积分数为
\[
\phi_b = \frac{1}{N} \sum_{\textbf x \in \Omega} B(\textbf x)
\]
其中 $N = |\Omega|$ 是计算域内的总体素数。给定目标体积分数 $\phi_{b,\textbf{target}}$，对应的目标 binder 体素数为
\[
N_b^{\star} = \phi_{b,\textbf{target}} \cdot N
\]
在生成的过程中，本文已写入的 binder 体素数
\[
N_b = \sum_{\textbf x \in \Omega} B(\textbf x)
\]
并据此更新当前体积分数 $\phi_b = N_b / N$。

当 $N_b \ge N_b^{\star}$ 时，认为已达到或超过目标体积分数，此时停止生成过程，也即
\[
|\phi_b - \phi_{b,\textbf{target}}| \le \delta_b
\]
成立时，停止 binder 生成并输出 $B(\textbf x)$。其中 $\delta_b$ 是一个小的容差参数，用于控制体积分数匹配的精度。为避免无限迭代，算法同时设置候选接触对遍历上限（或最大迭代次数）；当候选接触对集合 $\mathcal{C}$ 被遍历完成仍未达到目标时，报告当前可实现的 $\phi_b$ 并终止。

由于 binder 以离散体素写入，$N_b$ 可能出现超过 $N_b^{\star}$ 的情况。此时本文对多余的 binder 体素进行裁剪，使 $N_b$ 回到目标值（例如在不破坏局部连通性的前提下随机删除或按优先级删除）。
若 $N_b < N_b^{\star}$ 且候选接触对已用尽，则说明在当前纤维骨架与接触阈值设定下可生成的 binder 体积受限；此时可选择（i）保持当前 $B(x)$ 并报告误差，或（ii）采用轻微的补量策略（如沿纤维表面局部涂覆）以逼近目标。本文采用的具体策略在实现中给定，并在结果中记录最终误差。

\subsubsection{\textbf{输出字段与可视化变量}}

为便于后续仿真与可视化，本文在输出数据中同时提供离散相状态函数与用于平滑显示的连续场。具体地，纤维相与 binder 相分别由二值体素场 $S(\textbf x)$ 和 $B(\textbf x)$ 表示，孔隙相由 $P(\textbf x) = 1 - S(\textbf x) - B(\textbf x)$ 表示。上述二值场用于体积分数/孔隙率激素按及作为后续数值模拟的相区分依据。

与此同时，为了降低体素边界带来的锯齿感并获得更平滑的等值面显示，本文额外构造用于可视化的连续场 $\texttt{Fiber}$ 与 $\texttt{Binder}$，取值范围为 $[0,1]$，其可由二值相场经局部平滑/软边界映射得到（具体详见 2.4 节）。并定义合成显示场
\[
\texttt{Flag}=\max(\texttt{Fiber},\texttt{Binder})
\]
用于整体结构的快速预览。需要指出的是，$\texttt{Fiber}$ 和 $\texttt{Binder}$ 仅用于可视化展示，不改变 $S(\textbf x)$ 和 $B(\textbf x)$ 的体积分数统计结果。

在 Tecplot 中，为分别显示纤维与 binder，建议对 \texttt{Fiber} 与 \texttt{Binder} 分别绘制等值面（Iso-surface），并采用统一阈值 $\mathrm{Iso}=0.5$；同时可选用 \texttt{Allow Quads (Smoother Contours)} 与 Gouraud/Phong 光照以进一步减弱体素台阶带来的视觉伪影。

具体的平滑处理将会在 2.4 节中详细介绍，此处先行说明输出字段的定义与用途，以便读者理解后续的可视化结果。

\subsection{\textbf{平滑处理}}

由于纤维相/粘连剂在本文中以二值体素场 $S(\textbf{x})$ 和 $B(\textbf{x})$ 分别表示，直接绘制等值面会产生明显的锯齿状边界，无法满足 GDL 中 binder 的平滑连续特征。因此在不改变 2.1-2.3 的\textbf{物理统计与孔隙率计算}的前提下，我们将二值体素场转换为 $[0,1]$ 的连续可视化场 $\tilde{S}(\textbf{x})$ 和 $\tilde{B}(\textbf{x})$，并用以 Tecplot 展示。

\subsubsection{\textbf{平滑动机与适用对象}}

本文采用三维体素域 $\Omega \subset \mathbb{Z}^3$ 表征 GDL 微结构。对任意体素索引 $\textbf{x} = (i,j,k) \in \Omega$，纤维相与粘连剂分别由二值状态函数 $S(\textbf{x}),B(\textbf{x}) \in \{0,1\}$ 表示。孔隙相由 $P(\textbf{x}) = 1 - S(\textbf{x}) - B(\textbf{x})$ 给出。如 2.1 节所定义，孔隙率与各相体积分数等统计量均基于上述二值相场直接计算；同时，$S(\textbf{x}),B(\textbf{x}),P(\textbf{x})$ 也作为后续数值模拟的相区分依据。

然而，若最直接对二值体素场提取等值面（例如在 Tecplot 中绘制 Iso-surface），固-孔界面将严格受体素网格约束，导致明显的”台阶效应/锯齿伪影”（staircase artifacts）。该伪影会影响三维结构的展示观感与局部细节辨识，尤其在纤维半径与体素尺寸同量级、或存在细小 binder 胶桥/团聚结构时更为突出。需要强调的是，这类锯齿来自离散表示与渲染方式，并不意味着几何生成或物理统计本身存在误差。

为降低体素边界带来的锯齿感并获得更加平滑的界面显示，本文在输出数据中额外构造用于可视化的连续标量场 $\tilde{S}(\textbf{x}), \tilde{B}(\textbf{x}) \in [0,1]$。其基本思想是：以最终生成的二值相场 $S(\textbf{x}), B(\textbf{x})$ 作为输入，通过局部平滑/软边界映射将“硬边界”转换为具有亚体素过渡带的“软边界”连续场，从而使得等值面在特定阈值（例如 $\mathrm{Iso}=0.5$）附近穿过原始边界位置，实现更稳定、平滑的界面重建。进一步地，为便于整体结构快速预览，定义合成显示场
\[
\texttt{Flag} = \max(\tilde{S}(\textbf{x}), \tilde{B}(\textbf{x}))
\]

本节平滑处理的\textbf{适用对象}为：\textbf{纤维相二值场 $S(\textbf{x})$} 与（若启用 binder）\textbf{粘连剂相二值场 $B(\textbf{x})$}。无论前述微结构生成采用“重叠/不重叠”模式，平滑步骤均在相场生成完成后对最终的 $S(\textbf{x})$ 和 $B(\textbf{x})$ 进行处理，因此不影响几何生成逻辑。尤为重要的是，$\tilde{S}(\textbf{x}),\tilde{B}(\textbf{x}),\texttt{Flag}$ 仅用于可视化展示，不参与孔隙率/体积分数等统计量的重新计量；本文所有的统计结果均以二值相场 $S(\textbf{x}), B(\textbf{x})$ 为基础计算，确保物理统计的准确性与一致性。

\subsubsection{\textbf{高斯平滑的数学形式}}

为降低体素离散导致的界面台阶伪影，并与 2.3.6 节的呃可视化字段（\texttt{Fiber/Binder/Flag}）保持一致，本文将二值指示场 $S(\textbf{x}), B(\textbf{x})$ 通过三维高斯平滑（Gaussian smoothing）进行平滑处理，映射为连续的可视化场 $\tilde{S}(\textbf{x}), \tilde{B}(\textbf{x})$。

三维高斯核定义为
\[
G_{\sigma}(\textbf{r}) = \frac{1}{(2\pi\sigma^2)^{3/2}} \exp\left(-\frac{\|\textbf{r}\|^2}{2\sigma^2}\right)
\]
其中 $\sigma$ 是标准差，控制平滑程度（单位为体素）；$\textbf{r} = (x,y,z)$ 是空间位置向量。对于离散体素场，平滑后的值通过卷积计算得到：该核满足
\[
\int_{\mathbb{R}^3} G_{\sigma}(\textbf{r}) \, d\textbf{r} = 1
\]
从而保证平滑操作作为局部加权平均，不引入整体强度偏移。

记 $F_0(\textbf x) = S(\textbf x)$、$H_0(\textbf x) = B(\textbf x)$，则平滑后的连续场 $\tilde{S}(\textbf{x}), \tilde{B}(\textbf{x})$ 通过以下卷积计算得到：
\[
\begin{cases}
\tilde{S}(\textbf{x}) = (G_{\sigma} * F_0)(\textbf{x}) = \sum_{\textbf{r} \in \Omega} G_{\sigma}(\textbf{x}-\textbf{r}) S(\textbf{r})\\
\tilde{B}(\textbf{x}) = (G_{\sigma} * H_0)(\textbf{x}) = \sum_{\textbf{r} \in \Omega} G_{\sigma}(\textbf{x}-\textbf{r}) B(\textbf{r})
\end{cases}
\]
其中 $*$ 表示三维卷积。由于 $S,B$ 为二值场，平滑后在界面附近形成 $0$--$1$ 的连续过渡区，而在相内区域则趋近于 $1$，在孔隙区域则趋近于 $0$。通过调整 $\sigma$ 的大小，可以控制过渡带的宽度，从而实现不同程度的平滑效果。

本节只给出高斯平滑在该问题中的应用形式，有关三维高斯的具体算法推导详见附录 \ref{appendix:gaussian}。

\subsubsection{\textbf{数值实现细节}}

高斯平滑在体素网格 $\Omega \subset \mathbb{Z}^3$ 上实现。输入为最终生成的二值相场。为构造可视化连续场，我们首先将二值相场转为布尔掩码再转为浮点数
\[
\begin{cases}
M_f(\textbf{x}) = \mathbb{I}[S(\textbf{x})=1]\\
M_b(\textbf{x}) = \mathbb{I}[B(\textbf{x})=1]
\end{cases}
\]
其最终 $\mathbb{I}[\cdot]$ 是指示函数，输出为 $0$ 或 $1$。数值实现中采用 \texttt{float32} 以降低内存占用。

根据上文提到的卷积形式，本文在离散网格上使用三维高斯滤波算子近似实现卷积：
\[
\begin{cases}
F(\textbf{x}) \leftarrow \mathcal{G}_{\sigma}(M_f)(\textbf x)\\
H(\textbf{x}) \leftarrow \mathcal{G}_{\sigma}(M_b)(\textbf x)
\end{cases}
\]
其中 $\mathcal{G}_{\sigma}$ 表示对三维数组施加标准差为 $\sigma$ 的高斯滤波操作。该步骤将 0/1 硬边界映射为 $[0,1]$ 的连续软边界，从而在界面附近形成亚体素过渡带。

在进行高斯滤波的过程中，临近边界的点有可能会需要访问到体素网格边界之外的值。若简单的将边界外的值视为 0，可能会导致临近边界的固相平滑值杯拉低，从而使得界面位置“向内收缩”。为避免这种边界效应，本文在高斯滤波时采用邻边界扩展（nearest），将边界外的值设置为与边界内最近的值相同（即边界复制），从而保证平滑后的界面位置更接近原始二值场的位置。



\section{Settings}

\subsection{GDL 随机微结构生成方法与参数设置}

\subsubsection{微结构生成核心框架}
本文基于分层堆叠与各向异性扰动策略，开发了GDL碳纤维随机微结构生成算法，解决传统生成方法中“各向同性毛球效应”“体素硬边界锯齿化”“粘结剂（binder）散点分布”等问题。算法核心流程为：分层划分→纤维几何采样（方向+中心）→纤维体素化（支持重叠/不重叠模式）→binder 精准生成（交点胶团+毛细桥约束）→体素场平滑，最终输出0~1连续体素场（0为孔隙，1为固相），提升微结构真实性与可视化平滑度。

\subsubsection{基础几何与孔隙率参数（表2）}
微结构计算域尺寸设置为 \(I \times J \times K = 155 \times 300 \times 300\) 体素（对应物理尺度 µm 级，厚度轴自动识别为最小维度，即 \(I\) 轴），总孔隙数 \(N=I \times J \times K\)。目标总孔隙率（含纤维与binder）设定为 \(\varepsilon_{\text{target,total}} = 0.8018\)，孔隙率收敛精度 \(\Delta\varepsilon = 0.0005\)，避免生成过程中孔隙率波动过大。

碳纤维几何参数参考实际GDL碳纤维特性，具体设置如表\ref{tab:basic_params} 所示。

\begin{table}[H]
\centering
\small
\begin{adjustbox}{width=\columnwidth}
\begin{tabular}{lll}
\toprule
\textbf{符号/名称} & \textbf{含义} & \textbf{典型取值/范围} \\
\midrule
\(I,J,K\) & 体素域尺寸 & \(155 \times 300 \times 300\) \\
\(r_{\text{fiber}}\) & 纤维半径（体素） & 3.5 \\
\(L_{\text{min}}, L_{\text{max}}\) & 纤维长度范围 & 140.0，320.0 \\
\texttt{fiber\_end\_mode} & 端点处理方式 & \texttt{clipped} \\
\texttt{max\_fibers} & 纤维生成上限 & 50000 根 \\
\(\varepsilon_{\text{target,total}}\) & 目标总孔隙率 & 0.8018 \\
\(\Delta\varepsilon\) & 孔隙率收敛精度 & 0.0005 \\
\bottomrule
\end{tabular}
\end{adjustbox}
\caption{基础几何与孔隙率参数}
\label{tab:basic_params}
\end{table}

{\raggedright
\hspace{2em}纤维长度采用长纤维裁剪模式（\texttt{clipped}），该模式先生成长纤维中心线再裁剪至计算域，使纤维端点主要分布在计算域边界，更贴近真实GDL子区域截取特征。
\par}
% 修正后的表4：自适应列宽 + 限制总宽度 + 自动换行
\subsection{分层与各向异性参数（表3）}
为模拟GDL“分层堆叠”的宏观特征，沿厚度轴（$I$ 轴）划分分层数 $n_{\text{layers}}$，分层厚度 $h_{\text{layer}} = \max(4.0 \times r_{\text{fiber}}, 1.0)$ 体素，最终分层数取 $\max(3, \text{round}(\text{thickness}/h_{\text{layer}}))$（至少3层保证分层效果）。层内纤维中心在厚度轴方向的扰动（jitter）设为 $0.35 \times h_{\text{layer}}$，平衡分层规整性与真实随机性。

各向异性扰动参数控制纤维方向的“乱中有序”，具体设置如表 \ref{tab:anisotropic_params} 所示。
\begin{table}[H]
\centering
\Huge
\begin{adjustbox}{width=\columnwidth}
\begin{tabular}{lll}
\toprule
\textbf{符号/名称} & \textbf{含义} & \textbf{取值/范围} \\
\midrule
$n_{\text{layers}}$ & 分层数 & $\max(3, \text{round}(\text{thickness}/h_{\text{layer}}))$ \\
$\sigma_z$ (jitter) & 层内厚度方向抖动 & $0.35 \times h_{\text{layer}}$ \\
$\sigma_{\alpha}$ & 面内角度离散度 & $55.0^\circ$ \\
$\sigma_{\beta}, \beta_{\max}$ & 离层倾角控制 & $\sigma_{\beta}=4.0^\circ,\ \beta_{\max}=10.0^\circ$ \\
相邻层方向步长 & 层间方向变化步长 & $20^\circ \text{（随机游走）}$ \\
\bottomrule
\end{tabular}
\end{adjustbox}
\caption{分层与各向异性参数}
\label{tab:anisotropic_params}
\end{table}

{\raggedright
\hspace{2em}面内角度扰动标准差\(\sigma_{\alpha}=55.0^\circ\)，该值越大面内方向越随机，当取值为 \(180^\circ\) 时近似各向同性；离层面倾角扰动标准差 \(\sigma_{\beta}=4.0^\circ\)，倾角硬截断阈值 \(\beta_{\text{max}}=10.0^\circ\)，可限制纤维沿厚度轴的分量，保证纤维主要平行于层面。每层平均面内方向采用“随机游走”生成，相邻层平均方向步长 \(20^\circ\)，使层间方向缓慢变化，强化分层视觉特征。
\par}


\subsection{binder 生成核心参数（表4）}
{\raggedright
\hspace{2em}pacbinder 体积分数（相对总体）设定为 \(\phi_{\text{binder,total}} = 1.2\%\)（0.012），参考 GDL 中 binder 占固相体积的典型比例，具体参数如表 \ref{tab:binder_params} 所示。
\par}
\begin{table}[H]
\centering
\small
\begin{adjustbox}{width=\columnwidth}
\begin{tabular}{lll}
\toprule a
\textbf{符号/名称} & \textbf{含义} & \textbf{取值/范围} \\
\midrule
\(\phi_{\text{binder,total}}\) & binder体积分数 & 0.012 \\
\(g_{\text{contact}}\) & 接触间隙 & \(0.35 \times r_{\text{fiber}}\) \\
\(r_{\text{binder,blob}}\) & 初始胶团半径 & \(1.05 \times r_{\text{fiber}}\) \\
\(d_{\text{user}}\) & 毛细桥连线长度 & \(2.65 \times r_{\text{fiber}}\) \\
\(\text{d}t, \text{d}s\) & 采样密度（体素） & 0.65，0.75 \\
约束条件 & 体素范围约束 & 纤维间隙slab区域 \\
\bottomrule
\end{tabular}
\end{adjustbox}
\caption{binder 生成核心参数}
\label{tab:binder_params}
\end{table}

{\raggedright
\hspace{2em}binder 生成时，仅在两根纤维“夹缝区”生成，体素需同时满足到两根纤维中心线距离 \(\leq \text{thr}\)（\(\text{thr}^2\) 为距离平方阈值），且限制在纤维间隙的slab区域内（投影范围 \([-{\text{margin}}, g+{\text{margin}}]\)），避免binder在单纤维表面形成凸起，保证胶桥的凹形圆角特征。
\par}

\subsection{数值计算与优化参数（表5）}
为提升算法效率与结果可复现性，采用多项数值优化策略，具体设置如表 \ref{tab:optimization_params} 所示。
\\
% 修正后的表5：同样适配单栏宽度
\begin{table}[H]
\centering
\begin{tabular}{p{1.8cm}p{2.7cm}p{2.5cm}}
\toprule
策略名称 & 含义 & 设置方式 \\
\midrule
Numba JIT 编译 & 核心几何计算加速 & 加速 point segment dist2 等函数 \\
AABB 分箱预筛 + 微调 & 不重叠模式优化 & 降低高密度碰撞检测开销 \\
Scipy 距离变换 & 体素平滑方式 & 生成窄过渡带（无则降级） \\
全局随机种子 & 可复现性控制 & 42 \\
\bottomrule
\end{tabular}
\caption{数值计算与优化参数}
\label{tab:optimization_params}
\end{table}

采用Numba即时编译（JIT）加速点-线段距离平方（\texttt{point\_segment\_dist2}）、线段-线段最近距离（\texttt{\_segseg\_closest}）等核心几何计算函数，可显著提升体素化效率；不重叠模式下采用轴对齐包围盒（AABB）分箱预筛+厚度轴微调（nudge）策略，能有效降低高密度下纤维碰撞检测的计算开销；体素平滑依赖Scipy的距离变换（\texttt{ndimage}）生成窄过渡带，若无Scipy则降级为无平滑模式，输出连续体素场替代传统0/1/2硬边界，提升Tecplot等值面可视化的光滑度；全局随机种子设为42，保证生成结果的可复现性。

\urlstyle{same}
\subsection{输出设置（表6）}
{\raggedright
\hspace{2em}算法生成4类微结构文件（\texttt{gdl\_layered\-\_opt\-\{opt\}.dat}），覆盖“重叠/不重叠”与“有/无 binder”的组合，具体如表 \ref{tab:output_files} 所示。
\par}

\begin{table}[H]
\centering
\begin{adjustbox}{width=\columnwidth}
\Large
\begin{tabular}{lll}
\toprule
\textbf{文件编号} & \textbf{文件名} & \textbf{组合模式} \\
\midrule
opt=1 & \path{gdl_layered_opt1.dat} & 不重叠+无 binder \\
opt=2 & \path{gdl_layered_opt2.dat} & 可重叠+无 binder \\
opt=3 & \path{gdl_layered_opt3.dat} & 不重叠+有 binder \\
opt=4 & \path{gdl_layered_opt4.dat} & 可重叠+有 binder \\
\bottomrule
\end{tabular}
\end{adjustbox}
\caption{输出文件设置}
\label{tab:output_files}
\end{table}
{\raggedright
\hspace{2em}所有文件均输出0~1连续体素场，可直接导入Tecplot、ParaView等可视化软件进行后处理，或用于孔隙尺度流动/传质模拟。
\par}

\section{Results}

%\input{results.tex}

\section{Discussion}
本次研究构建了支持\textbf{重叠/不重叠纤维生成}与\textbf{薄膜桥状binder建模}的GDL三维微结构生成方法，完成四种模式的微结构生成且孔隙率均逼近目标值，验证了方法的可行性，同时明确了当前模型存在的局限性与后续优化方向。

\subsection{研究的改进与创新特征}
本研究在微结构建模方法上实现多项改进与创新：提出厚度分层与高斯抖动采样法，突破碳棒纤维随机分布假设，精准复现 GDL 层状堆叠的结构特征；融合空间分箱、AABB 预筛与参数自适应收缩策略，有效解决不重叠模式下的 jamming 问题，提升模型生成稳定性；建立纤维接触几何驱动的 binder 靶向建模法，摒弃传统均匀填充思路，实现胶膜与胶桥的差异化形态构建；设计二值体素场软边界平滑算法，缓解体素化带来的锯齿伪影，显著优化微结构可视化效果；开发多模式切换建模框架并制定标准化输出流程，突破单一建模策略局限，有效提升方法的通用性与复用性。

\subsection{研究存在的局限性}
本研究仍存在若干不足之处。binder 形态表征的真实性有待提升，不重叠模式下生成的 binder 多以异面胶桥为主，而在重叠模式下局部区域易出现凸起特征，与实际结构分布存在偏差；不重叠模式（opt1/opt3）仍存在多次采样失败与 jamming 问题，算法整体生成效率仍有提升空间。现有平滑算法对局部细节的处理能力有限，在细小 binder 胶桥、纤维端部等区域仍存在较为明显的锯齿伪影；不重叠模式下生成的碳纤维中短细纤维占比偏高，且整体分布偏规整，与实际 GDL 碳纤维随机堆叠的特征存在差异。同时，模型将碳纤维简化为等半径理想长直圆柱体，未考虑纤维弯曲、粗细不均等真实几何特征，一定程度上影响了接触几何表征的准确性。此外，当前建模框架尚未引入 binder 亲疏水性、热导率等关键物理参数，暂难以实现 GDL 结构–物理–性能的多场耦合分析。

\subsection{后续优化方向}
后续将围绕微结构建模与仿真开展系统性优化：通过优化 binder 形态分类规则与厚度函数参数、提升纤维共面接触对占比并强化凹弧控制以修正局部凸起问题；引入 KD 树空间索引与多参数分级收缩策略，降低非重叠模式下碰撞检测计算开销并减少 jamming 现象；采用几何轮廓引导的自适应平滑算法，对局部细节区域精细化调控平滑核以提升界面平滑度；构建纤维半径分布与弯曲几何模型并引入随机粗糙度扰动，使碳纤维表征更贴近真实结构特征；建立 binder 物理性质赋值机制，明确几何形态与物理参数的关联关系，夯实多物理场建模基础；同时结合实验表征数据校准模型输入参数，通过定量指标验证微结构真实性，最终形成实验与仿真相互支撑、迭代完善的闭环优化体系。

\subsection{结论}
本次研究通过厚度分层采样、碰撞检测加速、binder靶向生成等策略，实现了GDL纤维相、binder相、孔隙相的三相微结构耦合表征，在层状结构复现、binder形态精准建模等方面的改进具有实际应用意义。但本方法在形态真实性、算法生成效率、物理性质表征等方面仍存在局限性，后续将通过模型精细化、算法优化与实验数据校准进一步提升方法性能，为GDL结构-性能耦合分析提供更可靠的几何模型支撑。
%本次研究构建了支持\textbf{重叠/不重叠纤维生成}与\textbf{薄膜桥状binder建模}的GDL三维微结构生成方法，完成四种模式的微结构生成且孔隙率均逼近目标值，验证了方法的可行性，同时明确了当前模型存在的局限性与后续优化方向。

\subsection{研究的改进与创新特征}
本研究在微结构建模方法上实现多项改进与创新：提出厚度分层与高斯抖动采样法，突破碳棒纤维随机分布假设，精准复现 GDL 层状堆叠的结构特征；融合空间分箱、AABB 预筛与参数自适应收缩策略，有效解决不重叠模式下的 jamming 问题，提升模型生成稳定性；建立纤维接触几何驱动的 binder 靶向建模法，摒弃传统均匀填充思路，实现胶膜与胶桥的差异化形态构建；设计二值体素场软边界平滑算法，缓解体素化带来的锯齿伪影，显著优化微结构可视化效果；开发多模式切换建模框架并制定标准化输出流程，突破单一建模策略局限，有效提升方法的通用性与复用性。

\subsection{研究存在的局限性}
本研究仍存在若干不足之处。binder 形态表征的真实性有待提升，不重叠模式下生成的 binder 多以异面胶桥为主，而在重叠模式下局部区域易出现凸起特征，与实际结构分布存在偏差；不重叠模式（opt1/opt3）仍存在多次采样失败与 jamming 问题，算法整体生成效率仍有提升空间。现有平滑算法对局部细节的处理能力有限，在细小 binder 胶桥、纤维端部等区域仍存在较为明显的锯齿伪影；不重叠模式下生成的碳纤维中短细纤维占比偏高，且整体分布偏规整，与实际 GDL 碳纤维随机堆叠的特征存在差异。同时，模型将碳纤维简化为等半径理想长直圆柱体，未考虑纤维弯曲、粗细不均等真实几何特征，一定程度上影响了接触几何表征的准确性。此外，当前建模框架尚未引入 binder 亲疏水性、热导率等关键物理参数，暂难以实现 GDL 结构–物理–性能的多场耦合分析。

\subsection{后续优化方向}
后续将围绕微结构建模与仿真开展系统性优化：通过优化 binder 形态分类规则与厚度函数参数、提升纤维共面接触对占比并强化凹弧控制以修正局部凸起问题；引入 KD 树空间索引与多参数分级收缩策略，降低非重叠模式下碰撞检测计算开销并减少 jamming 现象；采用几何轮廓引导的自适应平滑算法，对局部细节区域精细化调控平滑核以提升界面平滑度；构建纤维半径分布与弯曲几何模型并引入随机粗糙度扰动，使碳纤维表征更贴近真实结构特征；建立 binder 物理性质赋值机制，明确几何形态与物理参数的关联关系，夯实多物理场建模基础；同时结合实验表征数据校准模型输入参数，通过定量指标验证微结构真实性，最终形成实验与仿真相互支撑、迭代完善的闭环优化体系。

\subsection{结论}
本次研究通过厚度分层采样、碰撞检测加速、binder靶向生成等策略，实现了GDL纤维相、binder相、孔隙相的三相微结构耦合表征，在层状结构复现、binder形态精准建模等方面的改进具有实际应用意义。但本方法在形态真实性、算法生成效率、物理性质表征等方面仍存在局限性，后续将通过模型精细化、算法优化与实验数据校准进一步提升方法性能，为GDL结构-性能耦合分析提供更可靠的几何模型支撑。

\section{Conclusion}

%\input{conclusion.tex}

\section{}


%%%%%% Figures %%%%%%
% \subsection*{Figures}

% Figures should be called out within the text and numbered in the order of their citation in the text. Every figure must have a descriptive title beginning with ``Figure [Number] …'' All figure titles should be either a phrase or a sentence; do not mix the two styles. See Figure \ref{fig:1} for example.

% % Single-column figure (uses one column)
% \begin{figure}[t]
%     \centering
%     \includegraphics[width=0.45\textwidth]{fig 1}
%     \caption{This is an example single-column figure.}
%     \label{fig:1}
% \end{figure}

% Figures should be displayed on a white background. When preparing figures, consider that they can occupy either a single column (half page width) or two columns (full page width), and should be sized accordingly.

% If a figure consists of multiple panels, they should be ordered logically and labelled with roman letters (i.e., A, B, C, etc.). All labels should be explained in the legend. See Figure \ref{fig:2} for example.

% % Two-column-wide figure (spans both columns)
% \begin{figure*}[t]
%     \centering
%     \begin{subfigure}{0.4\textwidth}
%         \includegraphics[width=\textwidth, height=2in]{fig 1}
%         \caption{}
%         \label{fig:2a}
%     \end{subfigure}
%     \hfill
%     \begin{subfigure}{0.4\textwidth}
%         \includegraphics[width=\textwidth, height=2in]{fig 2}
%         \caption{}
%         \label{fig:2b}
%     \end{subfigure}
%     \caption{This is an example of a two-column-wide figure consisting of multiple panels. (\subref{fig:2a}) This is the first panel. (\subref{fig:2b}) This is the second panel.}
%     \label{fig:2}
% \end{figure*}

% Upon acceptance, authors will be asked to provide the figures as separate electronic files. A figure guide is available for reference from the \href{https://spj.science.org/pb-assets/SPJ/CustomPages/Misc/SPJ_Figure_Preparation_Guide}{Science Partner Journals website}. In summary, authors must supply figures as Adobe Portable Document Format (PDF), Adobe Illustrator (AI), or Encapsulated PostScript (EPS) for illustrations or diagrams; Tagged Image File Format (TIFF), JPEG, PNG, PhotoShop (PSD), EPS, or PDF for photography or microscopy. Images should be of at least 300 dpi resolution, unless due to the limited resolution of a scientific instrument. If an image has labels, the image and labels should be embedded in separate layers.

% %%%%%% Tables %%%%%%
% \subsection*{Tables}
% Tables should supplement, not duplicate, the text. They should be called out consecutively within the text and numbered in the order of their citation in the text. 

% Every table must have a descriptive title beginning with ``Table [Number] …'' as noted in Table \ref{tab:1}. If numerical measurements are given, the units should be included in the column heading. Every vertical column should have a heading, followed by a unit of measure (if any) in parentheses. Units should not change within a column. Vertical rules should not be used. 

% Centered headings of the body of the table can be used to break the entries into groups. Do not use footnotes in column heads; include any such details in sentence form in the table legend. Footnotes should contain information relevant to specific cells of the table; use lowercase letters in alphabetical order, as needed: a, b, c, etc.

% % Two-column-wide table (spans both columns)
% \begin{table*}[t]
%     \caption{This is an example two-column-wide table.}    
%     \centering
%     \begin{tabular}{cccccc}
%             \toprule
%             Column 1 & Column 2 & Column 3 & Column 4 & Column 5 & Column 6 \\
%             \midrule
%             Cell 1 & Cell 2 & Cell 3 & Cell 4 & Cell 5 & Cell 6\\
%             Cell 7 & Cell 8 & Cell 9 & Cell 10 & Cell 11 & Cell 12 \\
%             \bottomrule
%     \end{tabular}
%     \label{tab:1}
% \end{table*}

% \section{Materials and Methods}
% The materials and methods section should provide sufficient information to allow replication of the results. This section should be broken up by subheadings. Under exceptional circumstances, when a particularly lengthy description is required, a portion of the materials and methods can be included in the Supplementary Materials. 

% \subsection{Experimental Design}
% Begin with a section titled Experimental Design describing the objectives and design of the study as well as prespecified components. 

% \subsection{Statistical Analysis}
% If applicable, include a section titled Statistical Analysis that fully describes the statistical methods with enough detail to enable a knowledgeable reader with access to the original data to verify the results. The values for N, P, and the specific statistical test performed for each experiment should be included in the appropriate figure legend or main text. 

% \subsection{Human and Animal Research}
% For investigations on humans, a statement must be including indicating that informed consent was obtained after the nature and possible consequences of the study was explained.

% For authors using experimental animals, a statement must be included indicating that the animals’ care was in accordance with institutional guidelines.

% \section{Results}
% The results should describe the experiments performed and the findings observed. The results section should be divided into subsections to delineate different experimental themes. 
% \begin{itemize}
%     \item All data should be presented in the Results. No data should be presented for the first time in the Discussion. Data (such as from Western blots) should be appropriately quantified.
%     \item Subheadings must be either all complete sentences or all phrases. They should be brief, ideally less than 10 words. Subheadings should not end in a period. Your paper may have as many subheadings as are necessary.
%     \item Figures and tables must be called out in numerical order. For example, the first mention of any panel of Fig. 3 cannot precede the first mention of all panels of Fig. 2. The supplementary figures (for example, fig. S1) and tables (table S1) must also be called out in numerical order. 
% \end{itemize}

% \section{Discussion}
% Include a Discussion that summarizes (but does not merely repeat) your conclusions and elaborates on their implications. There should be a paragraph outlining the limitations of your results and interpretation, as well as a discussion of the steps that need to be taken for the findings to be applied. Please avoid claims of priority. 

% \section*{Acknowledgments}
% Anyone who made a contribution to the research or manuscript, but who is not a listed author, should be acknowledged (with their permission). Types of acknowledgments include:

% \subsection*{General} 
% Thank others for any contributions, whether it be direct technical help or indirect assistance 

% \subsection*{Author Contributions} 
% Describe contributions of each author to the paper, using the first initial and full last name. 

% \medskip Examples:

% ``S. Zhang conceived the idea and designed the experiments.''

% ``E. F. Mustermann and J. F. Smith conducted the experiments.''

% ``All authors contributed equally to the writing of the manuscript.''

% \subsection*{Funding}
% Name financially supporting bodies (written out in full), followed by the funding awardee and associated grant numbers (if applicable) in square brackets. 

% \medskip Example: 

% ``This work was supported by the Engineering and Physical Sciences Research Council [grant numbers xxxx, yyyy]; the National Science Foundation [grant number zzzz]; and a Leverhulme Trust Research Project Grant.'' 

% \medskip
% If the research did not receive specific funding, but was performed as part of the employment of the authors, please name this employer. If the funder was involved in the manuscript writing, editing, approval, or decision to publish, please declare this.

% \subsection*{Conflicts of Interest}
% Conflicts of interest (COIs, also known as ``competing interests'') occur when issues outside research could be reasonably perceived to affect the neutrality or objectivity of the work or its assessment. 

% Authors must declare all potential interests – whether or not they actually had an influence – in a ‘Conflicts of Interest’ section, which should explain why the interest may be a conflict. Authors must declare current or recent funding (including for Article Processing Charges) and other payments, goods or services that might influence the work. All funding, whether a conflict or not, must be declared in a ``Funding Statement.'' The involvement of anyone other than the authors who 1) has an interest in the outcome of the work; 2) is affiliated to an organization with such an interest; or 3) was employed or paid by a funder, in the commissioning, conception, planning, design, conduct, or analysis of the work, the preparation or editing of the manuscript, or the decision to publish must be declared.

% If there are none, the authors should state ``The author(s) declare(s) that there is no conflict of interest regarding the publication of this article.'' Submitting authors are responsible for coauthors declaring their interests. Declared conflicts of interest will be considered by the editor and reviewers and included in the published article.

% \subsection*{Data Availability}
% A data availability statement is compulsory for all research articles. This statement describes whether and how others can access the data supporting the findings of the paper, including 1) what the nature of the data is, 2) where the data can be accessed, and 3) any restrictions on data access and why.

% If data are in an archive, include the accession number or a placeholder for it. Also include any materials that must be obtained through a Material Transfer Agreements (MTA). 

% \section*{Supplementary Materials}
% Describe any supplementary materials submitted with the manuscript (e.g., audio files, video clips or datasets). 

% Please group supplementary materials in the following order: materials and methods, figures, tables, and other files (such as movies, data, interactive images, or database files). 

% \medskip Example:
% Fig. S1. Title of the first supplementary figure.

% Fig. S2. Title of the second supplementary figure.

% Table S1. Title of the first supplementary table.

% Data file S1. Title of the first supplementary data file.

% Movie S1. Title of the first supplementary movie.

% \medskip
% Be sure to submit all supplementary materials with the manuscript and remember to reference the supplementary materials at appropriate points within the manuscript. We recommend citing specific items, rather than referring to the supplementary materials in general, for example: ``See Figures S1-S10 in the Supplementary Material for comprehensive image analysis.''

% A link to access the supplementary materials will be provided in the published article.

% Supplementary Materials may include additional author notes—for example, a list of group authors.

% \section*{Guidelines for References}
% Authors are responsible for ensuring that the information in each reference is complete and accurate. All data must be cited and references to ``data not shown'' or citations to unpublished results are permitted.

% All references should be cited within the text and uncited references will be removed.

% There is only one reference list for all sources cited in the main text, figure and table legends, and Supplementary Materials. Do not include a second reference list in the Supplementary Materials section. References cited only in the Supplementary Materials section are not counted toward length guidelines.

% Please do not include any extraneous language such as explanatory notes as part of a reference to a given source. The journal prefers that manuscripts do not include end notes; if information is important enough to include, please put into main text.  If you need to include notes, please explain why they are needed in your cover letter to the editor.

% DOIs, if available, should be included for each reference.

% \printbibliography

\end{document}
